\section{\enumo: A DSL for strategic theory exploration}
\label{sec:slide}

\renewcommand\denote[1]{\ensuremath{\left\llbracket #1 \right\rrbracket}}
\newcommand\set[2]{\ensuremath{\left\{#1 \hspace{3pt} \left\vert \hspace{3pt} #2 \right.\right\}}}
\newcommand\plugsexp{\ensuremath{\mathsf{plug\_sexp}}}

\begin{figure}
  \begin{subfigure}{\linewidth}
    \begin{minipage}[t]{0.45\linewidth}
      \begin{grammar}
      <workload> ::=  Set <s-exp>*
        \alt Union <workload>*
        \alt Filter <filter> <workload>
        \alt Plug <workload> <string> <workload>

      <s-exp> ::=  Atom <string> \alt List <s-exp>+

      \end{grammar}
    \end{minipage}\hfill%
    \begin{minipage}[t]{0.5\linewidth}
      \begin{grammar}

      <filter> ::=
        MetricLt <metric> <\N>
        \alt MetricEq <metric> <\N>
        \alt Contains <pattern>
        \alt Canon <string>+
        \alt And <filter>+ | Or <filter>+ | Not <filter>

      <metric> ::= Atoms | List | Depth
      \end{grammar}%
      \end{minipage}%
    \caption{\enumo workload abstract syntax.}
    \label{fig:workload-grammar}
  \end{subfigure}

  \begin{subfigure}{\linewidth}
    \begin{align*}
      \denote{ \mathsf{Set} ~ ts } &= ts
      &
      \denote{ \mathsf{Filter} ~ \mathit{filter} ~ \W } &=
      \set{t \in \denote{ \W }}{\denote{ \mathit{filter} } (t) = true}
      \\
      \denote{ \mathsf{Union} ~ \W_1 ~ \W_2 } &=
      \denote{ \W_1 }~\cup~\denote{ \W_2 } \hspace{8mm}
      &
      \denote{ \mathsf{Plug} ~ \W_1 ~ \mathit{tgt} ~ \W_2 } &=
        \bigcup_{e \in \denote{\W_1}}
        \denote{ \plugsexp ~ e ~ \mathit{tgt} ~ \W_2 }
    \end{align*}
    \begin{align*}
      \denote{ \plugsexp ~ (Atom ~ s) ~ \mathit{s} ~ \W } &=
      \denote{ \W }
      \\
      \denote{ \plugsexp ~ (Atom ~ s) ~ \mathit{tgt} ~ \W } &=
      \{Atom ~ s\} ~\mathsf{when} ~ s \neq \mathit{tgt}
      \\
      \denote{ \plugsexp ~ (List ~ s_1 ~ \ldots ~ s_n) ~ \mathit{tgt} ~ \W } &=
      \set{List ~ t_1 ~ \ldots ~ t_n}{t_i \in \denote{ \plugsexp ~ s_i ~ \mathit{tgt} ~ \W }}
    \end{align*}
    \caption{\enumo workload semantics.}
    \label{fig:denote-workload}
  \end{subfigure}

  \begin{subfigure}{\linewidth}
    \begin{align*}
      \denote{\mathsf{MetricLt} ~ M ~ n}(t) &= \denote{ M } (t) < n  &
      \denote{\mathsf{Contains} ~ p}(t) &=
        \exists \sigma, \sigma(p) \in \mathsf{subterms}(t) \\
      \denote{\mathsf{MetricEq} ~ M ~ n}(t) &= \denote{ M } (t) = n  &
      \denote{\mathsf{Canon} ~ \vec{a} ~ }(t) &= canon(\vec{a}, t) == t
    \end{align*}
    \caption{
      \enumo filter semantics.
      Metric filters test various measures of term size
        (number of atoms, number of lists, or depth), e.g.,
        \texttt{(+ a b)} has 3 atoms, 1 list, and depth 2.
      The $\mathsf{Contains}$ filter tests whether
        any subterm of $t$ matches a pattern $p$.
      The $\mathsf{Canon}$ filter tests whether
        $t$ is \textit{canonical} with respect to
        a vector of atoms $\vec{a}$, e.g.,
        $\mathsf{Canon} ~ [a, b, c] ~ \mathtt{(+ ~ a ~ b ~ c)}$ is true while
        $\mathsf{Canon} ~ [a, b, c] ~ \mathtt{(+ ~ b ~ a ~ c)}$ is false.
    }\label{fig:denote-filter}
  \end{subfigure}

  \caption{
    Syntax and semantics for the workload fragment of the \enumo DSL.
  }\label{fig:workload}
\end{figure}

This section presents the core of
  the \enumo DSL for
  guided term enumeration and
  incremental rewrite rule inference.
\enumo programs primarily manipulate two kinds of values:
  \textit{workloads} representing sets of terms and
  \textit{rulesets}, which are sets of pairs of patterns.
Both terms within workloads and
  patterns within rewrite rules are
  represented as (untyped) s-expressions.

At a high level, \enumo programs
  typically iterate the following steps:
\begin{enumerate}
  \item Construct a workload $\W$
    representing a search space
    with terms of interest.
  \item Convert $\W$ to an \egraph and
    use the current ruleset to
    merge equivalent \eclasses.
  \item Search this compressed \egraph
    to find candidate rewrite rules,
    i.e., unmerged pairs of \eclasses of terms
    that fuzzing or other techniques suggest
    may be equivalent.
  \item Minimize the candidates:
    remove rules that are redundant
    given the current ruleset.
  \item Add the set of minimized candidates to
    the current ruleset.
\end{enumerate}

\enumo provides a number of operators for
  constructing and manipulating both
  workloads and rulesets, including
  \textit{plugging} and \textit{iterating} workloads
    to build up sets of terms,
  \textit{forcing} to materialize and insert a
    workload of terms into an e-graph,
  \textit{searching} e-graphs built from workloads for
    candidate rewrite rules, and
  \textit{minimizing} rulesets to
    remove redundant rules.
\enumo programs then are just a sequence of bindings
 from variables to workload and ruleset expressions
 embedded in a host language,
 e.g., a simple lambda calculus.

\subsection{Workloads}
\label{subsec:workloads}

\autoref{fig:workload} shows the syntax and semantics of
  workloads in \enumo.
Workloads have four constructors:
  \textsf{Set} represents a literal set of \sexps,
  \textsf{Union} represents unions of workloads,
  \textsf{Filter} represents a subset of terms in a workload, and
  \textsf{Plug} represents substituting
    one workload into another.
Note that workloads only \textbf{represent} sets of terms:
  \C{force} can be applied to a workload to
  materialize the set of
  untyped \sexps (i.e., terms)
  it represents.
This laziness is a key design decision
  that allows \enumo programs to efficiently
  represent and manipulate large sets of terms.

\textsf{Set} and \textsf{Union}
  have straightforward semantics
  (\autoref{fig:denote-workload}).
\textsf{Filter}
 takes a workload $\W$ and a filter predicate $P$,
 and represents the set of terms from $\W$ that satisfy $P$.
\autoref{fig:workload-grammar} shows the
 syntax of filter predicates,
 some of which use term \textit{metrics}.
Metrics count
 the number of atoms, lists, and
 depth of a term.
The semantics for filters is given in \autoref{fig:denote-filter}.
The \textsf{MetricEq} and \textsf{MetricLt} filters
 measure a metric of a term and compare it to a given value.
The \textsf{Contains} filter
 checks whether the given pattern occurs in a term.
The \textsf{Canon} filter
 checks that a given term is canonical with respect
 to a given list of variables.
The \textsf{Not}, \textsf{Or}, and \textsf{And} filters
 are simply the usual logical connectives.

\textsf{Plug}
  allows the user to substitute all combinations
  of terms from one workload in for
  a given variable in another workload.
As the semantics in \autoref{fig:denote-workload} show,
  \textsf{Plug} provides a special kind of substitution
  that performs a Cartesian product:
  \textsf{Plug $\W_1$ $s$ $\W_2$}
  returns a workload
  that denotes a set of
  $\sum_{e \in \denote{\W_1}} |\denote{\W_2}|^{k_e}$ terms\footnote{
    Duplicate terms are removed, so this is an upper bound on the number
    of terms in the resulting workload.
  },
  where $k_e$ is the number of occurrences of $s$ in the term $e$.
The following \enumo snippet demonstrates
  the semantics of \textsf{Plug}:
\begin{lstlisting}
w1 = { X (foo X X) }
w2 = { 1 y }
plugged = w1.plug("X", w2) # { 1 y (foo 1 1) (foo 1 y) (foo y 1) (foo y y)}
\end{lstlisting}

\enumo's operators can be composed into useful,
  reusable strategies beyond the concise reimplementation
  of past work.
As an example,
  \T{iter_metric}, defined in the following \enumo snippet,
  can be used to create size-parameterized workloads:
  \T{iter_metric(W, tgt, Atoms, n)} produces all terms
  from the workload
  with at most $n$ atoms, and \T{iter_metric(W, tgt, Depth, n)}
  produces all terms from the workload up to depth $n$. \T{iter_metric}
  can be used to generate workloads with successively larger terms
  and thus guide the exploration of successively deeper rules
  across domains (\autoref{sec:eval}).

\begin{lstlisting}[
language=Python,
basicstyle=\footnotesize\ttfamily,
numbers=left,
xleftmargin=2.5em]
def iter_plug(W, tgt, n):
  if n <= 0: return W
  return W.plug(tgt, iter_plug(W, tgt, n - 1))

def iter_metric(W, tgt, metric, n):
  return iter_plug(W, tgt, n).filter(metric <= n)
\end{lstlisting}

\paragraph*{Optimizing workloads.}

\textsf{Plug} is the key workload combinator
  for representing search spaces by
  enumerating terms from a grammar.
\textsf{Plug} typically represents
  combinatorially many more terms
  than its arguments,
  but the result of a \textsf{Plug}
  is often \textsf{Filter}ed to
  target a more specific subset of
  the represented terms.
We introduce an essential optimization
  to speed up workload evaluation (forcing)
  that avoids unnecessary work during
  combinatorial substitution by
  pushing some \textsf{Filter}s through \textsf{Plug}s
  according to the following equation:
\begin{align*}
  \denote{\mathsf{Filter} ~ \mathit{filter} ~
          (\mathsf{Plug} ~ \W_1 ~ s ~ \W_2)} &=
  \denote{\mathsf{Filter} ~ \mathit{filter} ~
          (\mathsf{Plug} ~ \W_1 ~ s ~ (\mathsf{Filter} ~ \mathit{filter} ~ \W_2))}
\end{align*}
\noindent
when \textit{filter} is monotonic.
A filter $f$ is \emph{monotonic}
  if, for every term $t$ satisfying $f$,
  every subterm $s \in t$ also satisfies $f$.
Note that the outer \textsf{Filter}
  remains in place even after the optimization,
  as removing it entirely would not
  preserve semantics.
This still yields an exponential reduction
  in the number of terms that must be
  materialized and filtered.
\textit{
  All of the \enumo programs in our evaluation
  depend heavily on this optimization.}

In the current \enumo implementation,
  \textsf{And} and \textsf{MetricLt} filters (among others) are pushed
  through \textsf{Plug}s.
This optimization and its
  monotonicity constraint are inspired by the
  classic relational algebra optimization
  of pushing certain selections through joins~\citep{alice};
  \textsf{Plug}'s combinatorial behavior,
  in some ways, resembles a relational join.

\subsection{\Egraphs and rulesets}
\label{subsec:ruleset-sem}

\renewcommand{\to}{\textrightarrow\xspace}
\newcommand{\head}[1]{\textsf{\sc #1}}
\begin{figure}
  \footnotesize
  \begin{minipage}[t]{0.5\linewidth}
  \centering
    \raggedright\tt
    \head{E-graph Generating Operators} \\
    \vspace{4pt}
    \toeggl  : $\T{workload}$ \to \eggl \\
    $\T{eqsat}$    : \eggl \to \ruleset \to \eggl \\
    $\T{compress}$ : \eggl \to \ruleset \to \eggl \\[1mm]
    \vspace{4pt}
    \head{Rule Testing Operators} \\
    \vspace{4pt}
    $\T{is\_saturating}$ : \rl \to $\T{bool}$ \\
    $\T{is\_valid}$      : \rl \to $\T{bool}$ \\
    $\T{can\_derive}$    : \ruleset \to \rl \to $\T{bool}$
    \vspace{4pt}
  \end{minipage}
  \begin{minipage}[t]{0.65\linewidth}
  \centering
    \raggedright\tt
    \head{Ruleset Generating Operators} \\
    \vspace{4pt}
    $\T{find\_candidates}$    : \eggl \to \ruleset \\
    $\T{partition}$          : \ruleset \to (\rl \to $\T{bool}$) \to \ruleset $\ast$ \ruleset \\
    $\T{minimize}$           : \ruleset \to \ruleset \to \ruleset \\
    $\T{union}$              : \ruleset \to \ruleset \to \ruleset \\
    $\T{candidates\_by\_diff}$ : \eggl \to \eggl \to \ruleset
  \end{minipage}
  \caption{
    \enumo's operators over rulesets and \egraphs.
  }\label{fig:ruleops}
\end{figure}

In addition to novel, programmable term enumeration,
 \enumo also provides primitives to
 create and manipulate \egraphs and rulesets.
\autoref{fig:ruleops} shows
  these operators and their types.
Many mirror parts of
  earlier monolithic theory explorers;
\enumo's key insight lies in turning such tools
  ``inside out'' to provide their components
  as composable operators in a DSL that
  allows users to strategically guide
  the search for rewrites and
  incrementally build rulesets.

\paragraph*{\Egraph operators.}
For the purposes of the \enumo language definition,
  an \egraph is an abstract data type
  that provides the operations described
  in \autoref{sec:bg-egraphs}.
A typical \enumo program (\autoref{sec:overview})
 converts a workload into an \egraph
 using the \toeggl operator
 before proceeding to candidate generation.
The resulting \egraph represents every term in the set
 denoted by the workload.

The \texttt{eqsat} and \texttt{compress} operators
 both run equality saturation on the given \egraph
 with the given ruleset.
The former allows the \egraph to grow while the latter
  only applies merges over existing terms in the \egraph.
From \enumo's perspective,
 \texttt{compress}'s main purpose is to
 remove redundancy from the \egraph implied
 by a ruleset of already-learned rewrites.
Because we build on the \egg library,
 \texttt{eqsat} in our implementation
 is also parameterized over a \emph{strategy}
 that determines how a ruleset is applied to the \egraph
 (\autoref{subsec:cmp}).

\paragraph*{Ruleset operators.}

A ruleset is a set of rewrite rules,
 where each rule is a pair of patterns.
Rulesets can be
 read from or written to a file,
 manipulated using the ruleset operators in \autoref{fig:ruleops},
 and used to perform equality saturation on \egraphs
 using the \texttt{eqsat} operator described above.

In a typical \enumo program,
 the \T{find\_candidates} operator
 is used to infer a ruleset from an \egraph.
\T{find\_candidates}
 is parameterized on a user-provided interpreter
 which is used to
 identify likely sound rule candidates
 by evaluating the terms over a set of inputs (fuzzing);
 terms that disagree are certainly not equivalent,
 but those that agree may be~\citep{ruler,bansal}.
\renewcommand{\eval}{\textsf{eval}}
\newcommand{\repr}{\textsf{repr}}
To define \T{find\_candidates} formally,
 let $\repr(e)$ denote a \textit{representative} term from
 \eclass $e$
 and let $\eval(t)$ denote the result of evaluating $t$
 over a set of input values.
Then \T{find\_candidates}
  on \egraph $E$ returns a set of rules:
\[
  \set{\repr(e_l) \rewritesto \repr(e_r)}{
    \text{$e_l, e_r \in E$}.\
    \eval(\repr(e_l)) = \eval(\repr(e_r))
  }
\]

The \T{partition} operator takes a \ruleset $R$
  and a predicate $P$ over rules and
  returns $(R_1, R_2)$ such that
  $R = R_1 \cup R_2$,
  $r \in R_1 \rightarrow P(r)$, and
  $r \in R_2 \rightarrow \neg P(r)$.
\autoref{fig:ruleops}
 contains two such predicates, \T{is\_valid} and \T{is\_saturating}.
The \T{is\_valid} predicate
 checks whether a rule $\ell \rewritesto r$
 is valid for all inputs,
 using a user-provided verifier
 for the domain.
Depending on the domain,
 the verifier may use SMT, model checking,
 or other techniques like fuzzing.
The built-in \T{is\_saturating} predicate
 checks whether a rule is \textit{saturating},
 i.e., applying the rule to an \egraph
 will not increase the size of the \egraph,
 measured as the number of \eclasses.
At a high level,
  saturating rules are those whose right-hand side pattern
  only contains subterms that appear in the left-hand side,
  except potentially for the root operator.
For example,
  $x + y \rewritesto y + x$ is saturating
  since all non-root subterms in the right-hand side ($x$ and $y$) also
  occur in the left-hand side but
  $x + (y + z) \rewritesto (x + y) + z$ is not
  since the right-hand side contains a non-root subterm $x + y$
  that does not appear in the left-hand side.
Applying only saturating rules
  to an \egraph
  is guaranteed to reach a fixpoint
  past which further application of the rules
  will no longer change the \egraph.

The \T{can_derive} operator
 tests whether a \ruleset $R$
 can derive a rule $\ell \rewritesto r$,
 discussed below in \autoref{subsec:derivability}.
The \T{minimize} operator
 takes a ruleset $R$ and prior rules $P$
 and \textit{minimizes} $R$ with respect to $P$,
 guaranteeing that $
  \mathtt{minimize}(R, P) = R' \implies
  \forall r \in R \setminus R',
  \T{can\_derive}(P \cup R', r)
 $.
Conceptually,
 this filters $R$ to a small subset $R'$ of rules
 such that $r \in R' \rightarrow \neg \T{can\_derive}(P, r)$.
However, \enumo's \T{minimize} operator provides
 an optimization that batches these checks to simultaneously
 eliminate redundant rules, initially described in \citet{ruler}.

The final core operator
  is \T{candidates\_by\_diff},
  which takes two \egraphs $e_1$ and $e_2$ and
  returns an inferred ruleset.
\T{candidates\_by\_diff}
  infers candidate rewrite rules from
  \eclasses which merged during an \eqsat run,
  i.e., terms which a given ruleset could prove equivalent.
Typically, $e_2$ is the result of
  running \eqsat on $e_1$ with ruleset $R$.
If, by application of $R$, the equivalence between
  terms $t$ and $t'$ is discovered,
  then \T{candidates\_by\_diff} will learn a
  rule candidate by extracting the best expression from the \eclasses
  representing $t$ and $t'$ in $e_1$.
\T{candidates\_by\_diff} enables rule synthesis for new domains
  which prior work could not support. This utility is
  briefly exemplified in~\autoref{subsec:ffoverview},
  and formally presented in~\autoref{sec:lift}.

\subsection{Discussion on derivability}
\label{subsec:derivability}

Given two rulesets, how do we know which is better?
While it may be tempting to use ruleset size as a proxy
  for ruleset quality, more rules are not necessarily better
  because overly redundant rules lead to slower performance of \eqsat
  systems.
A small set of simple rules is often easier to maintain and
  debug than a large set of complicated rules.
On the other hand, a ruleset with too few rules is less useful
  because fewer equivalences will be found,
  especially since resource limits restrict
  the number of iterations of \eqsat.
Since saturation is rare in practice, it is often helpful to have
  \textit{some} redundancy in the rulesets to get better results
  under given resource limits (see \autoref{sec:lift}).
 A ruleset's \textit{proving power} under given resource
  limits is subtle and difficult to estimate.
In this section, we define ruleset \textit{derivability}, a
  metric for measuring proving power that we use for comparing rulesets.

\paragraph*{Derivability.}
Prior work has not established a standard definition of derivability
  in the context of \eqsat.
In this paper, we formalize two ``obvious'' definitions of derivability:
  \lhsandrhs and \lhs.
To test whether a ruleset $R$ can derive a rule $\ell \rewritesto r$
  under given resource limits, we use the \eqsat procedure
  (\autoref{fig:eqsat-alg}).
The function \C{timeout}
  determines when to stop the \eqsat loop based on available resources
  (e.g., node count, iteration, or time bounds).
If running \eqsat using ruleset $R$ causes \eclasses representing
  $\ell$ and $r$ to merge, we say $\ell \rewritesto r$ is \textit{derivable}
  from $R$ under given resource bounds.
The \lhsandrhs derivability metric measures whether the equivalence between
  $\ell$ and $r$ can be recovered by applying the rules in $R$ to
  an \egraph initialized with both $\ell$ and $r$.
In contrast, the stronger \lhs definition for derivability states
  that $\ell \rewritesto r$ can be recovered given only $\ell$.
Prior work used \lhsandrhs derivability~\citep{ruler}.

In the context of \eqsat, the initial state of the \egraph interacts with
  resource limits in subtle ways because it changes what terms are
  available during e-matching.
Rules in $R$ must find concrete terms in the \egraph that match the
  left side of the rule in order to add the right side and merge the two
  \eclasses.
Changing the initialization of the \egraph changes what rule matches
  are possible.

To illustrate the difference between \lhs and \lhsandrhs derivability,
  consider the rule $a \rewritesto b$
  (where $a$ and $b$ are arbitrary patterns) and a ruleset containing
  the rule $b \rewritesto a$.
In an \egraph initialized with both $a$ and $b$ (\lhsandrhs), the rule
  $b \rewritesto a$ will fire and the \eclasses will merge, so the
  rule $a \rewritesto b$ will be considered derivable.
In an \egraph initialized with just $a$ (\lhs), $b \rewritesto a$
  will not fire, so the rule $a \rewritesto b$ will not be considered
  derivable.
\lhs and \lhsandrhs derivabilities may also require different resource
  limits.
For example, consider using
  $R = \{a \rewritesto b, b \rewritesto c, c \rewritesto b \}$ to
  derive the rule $a \rewritesto c$.
Under \lhsandrhs, the \eclasses representing $a$ and $c$ will
  merge within a single iteration of \eqsat.
In contrast, using \lhs derivability, recovering the equality between $a$ and
  $c$ will take two iterations of \eqsat.
First, the rule $a \rewritesto b$ will fire, creating an \eclass for
  $b$ and merging it with $a$'s \eclass.
In the second iteration, the rule $b \rewritesto c$ will fire, creating
  an \eclass for $c$ and merging it with the \eclass that represents
  $a$ and $b$, thus recovering the equivalence between $a$ and $c$.
This example shows that \lhsandrhs derivability may be able to derive equivalences
  in fewer iterations (i.e., using fewer resources) than \lhs because it can
  match on the left-hand and right-hand sides simultaneously.

In general, \lhs derivability is more conservative.
Anecdotally, we find that it is preferable
  when the user is interested in optimization-based
  \eqsat applications where an \egraph is initialized
  with a single term $t$, and \eqsat is used to find
  a better, equivalent version of $t$.
In contrast, \lhsandrhs derivability is looser,
  but may be appropriate in equivalence-checking
  \eqsat applications where two terms $t_1$ and $t_2$
  are added to an \egraph and \eqsat is used only to
  determine whether their \eclasses merge.

